% Chapter 12: Multilevel Monte Carlo.

\Cref{ch:07-path-simulation} ended with an uncomfortable trade-off. When the solution of a
stochastic differential equation cannot be sampled exactly, a discretised Monte Carlo
estimator pays twice: the bias of the scheme calls for small time steps, the statistical
error for many paths, and with a first-order scheme the cost of a root-mean-square error
$\varepsilon$ grows like $\varepsilon^{-3}$ instead of the $\varepsilon^{-2}$ of exact
sampling (\cref{prop:ps-allocation,prop:mc-mse}). Multilevel Monte Carlo closes most of this
gap, and in favourable cases all of it.

The idea fits in one line. Write the price computed with the finest step as the price
computed with a very coarse step plus a sum of corrections, each the difference between the
payoffs computed with two consecutive steps \emph{on the same Brownian path}. The coarse
price is cheap, so it can be estimated from many paths. A correction is expensive per path,
but the two paths in it are close, so its variance is small and few paths suffice. Most of
the work is done where it is cheap.

\Citet{giles-mlmc} introduced the method for path simulation, after the multilevel
methods of \textcite{heinrich-mlmc} for parametric integration and the two-level
``statistical Romberg'' estimator of \textcite{kebaier-romberg}. We prove the identity it
rests on and the best allocation of samples across levels (\cref{sec:ml-identity}), measure
how fast the corrections shrink under the Euler and Milstein couplings, also for
discontinuous payoffs (\cref{sec:ml-coupling}), prove Giles' complexity theorem with its
three regimes and compare it with the allocation of \cref{prop:ps-allocation}
(\cref{sec:ml-complexity}), describe the adaptive algorithm used in practice
(\cref{sec:ml-algorithm}), and end with the randomised estimators of
\textcite{rhee-glynn}, which remove the bias altogether when the corrections shrink fast
enough (\cref{sec:ml-unbiased}).

Throughout, $X$ solves \eqref{eq:ps-sde} on $[0,T]$ under $\Q$ with the Brownian motion $W$;
$P$ is a square-integrable discounted payoff, a function of $X_T$ or of the path of $X$;
$\theta=\E[P]$ is its price; and $P_\ell$ is the payoff computed from a scheme with
step $h_\ell=T2^{-\ell}$, $\ell\in\N$. The running example is the test case: geometric
Brownian motion $\dd S=\mu S\,\dd t+\sigma S\,\dd W$ with $\mu=r-y=-1\%$, $\sigma=20\%$,
$T=1$, and $P=e^{-rT}(S_T-K)^+$ with $S_0=K=100$, so $\theta=7.36429$. Costs are counted in
time steps: one step of one path costs $1$. A warning on notation: \cref{def:ps-orders}
called the strong and the weak orders $\gamma$ and $\beta$. From now on, as in the
literature on multilevel methods, $\alpha$, $\beta$ and $\gamma$ are the rates at which the
bias, the variance of the corrections and the cost change from one level to the next, and
the strong order of a scheme is written $\lambda$.

\section{The multilevel identity}\label{sec:ml-identity}

For a finest level $L$, linearity of expectation gives
\begin{equation}\label{eq:ml-telescoping}
  \E[P_L]=\E[P_0]+\sum_{\ell=1}^L\E\big[P_\ell-P_{\ell-1}\big].
\end{equation}
This is a triviality. It becomes a method when every term on the right is estimated from
its own independent samples, and every difference is computed from a fine and a coarse
approximation that are strongly dependent.

\begin{definition}[Coupled corrections]\label{def:ml-coupled}
Let $P_0,P_1,\dots\in L^2$. A \emph{coupled correction} at level $\ell\ge1$ is a random
variable $Y_\ell=P^{\mathrm f}_\ell-P^{\mathrm c}_{\ell-1}$, where the \emph{fine} member
$P^{\mathrm f}_\ell$ and the \emph{coarse} member $P^{\mathrm c}_{\ell-1}$ are square-integrable
random variables on one probability space with
\[
  \E\big[P^{\mathrm f}_\ell\big]=\E[P_\ell],\qquad \E\big[P^{\mathrm c}_{\ell-1}\big]=\E[P_{\ell-1}];
\]
usually $P^{\mathrm f}_\ell$ has the law of $P_\ell$ and $P^{\mathrm c}_{\ell-1}$ the law of
$P_{\ell-1}$. At level $0$, $Y_0=P_0$. We write $V_\ell=\Var(Y_\ell)$ and $\tau_\ell>0$ for
the expected cost of one sample of $Y_\ell$.
\end{definition}

Only the means of the two members are constrained; their joint law is ours to choose, and
the whole art is to choose it so that $V_\ell$ is small.

\begin{proposition}[The multilevel estimator]\label{prop:ml-telescoping}
Let $Y_0,Y_1,\dots$ be coupled corrections, $L\in\N$ and $n_0,\dots,n_L\ge1$ integers. Let
$Y^{(i)}_\ell$, $1\le i\le n_\ell$, $0\le\ell\le L$, be independent, with $Y^{(i)}_\ell$
distributed as $Y_\ell$, and
\begin{equation}\label{eq:ml-estimator}
  \hat\theta_L=\sum_{\ell=0}^L\hat Y_\ell,\qquad \hat Y_\ell=\frac1{n_\ell}\sum_{i=1}^{n_\ell}Y^{(i)}_\ell .
\end{equation}
\begin{enumerate}[label=(\roman*)]
  \item $\E[\hat\theta_L]=\E[P_L]$, whatever the joint laws of the fine and coarse members.
  \item $\Var(\hat\theta_L)=\sum_{\ell=0}^LV_\ell/n_\ell$, and for $\theta=\E[P]$ with
    $P\in L^1$,
    $\E\big[(\hat\theta_L-\theta)^2\big]=\big(\E[P_L]-\theta\big)^2+\sum_{\ell=0}^LV_\ell/n_\ell$.
  \item The expected cost is $\mathcal C=\sum_{\ell=0}^Ln_\ell\tau_\ell$. For all real
    $n_\ell>0$,
    \[
      \Big(\sum_{\ell=0}^L\frac{V_\ell}{n_\ell}\Big)\Big(\sum_{\ell=0}^Ln_\ell\tau_\ell\Big)\ \ge\ S_L^2,
      \qquad S_L=\sum_{\ell=0}^L\sqrt{V_\ell\tau_\ell},
    \]
    and if every $V_\ell>0$, equality holds if and only if $n_\ell$ is proportional to
    $\sqrt{V_\ell/\tau_\ell}$.
  \item Let $v>0$. Every choice of the $n_\ell$ with $\Var(\hat\theta_L)\le v$ costs at least
    $S_L^2/v$, and the integers
    \begin{equation}\label{eq:ml-nl}
      n_\ell=\max\Big(1,\Big\lceil\frac{S_L}{v}\sqrt{\frac{V_\ell}{\tau_\ell}}\,\Big\rceil\Big)
    \end{equation}
    give $\Var(\hat\theta_L)\le v$ at a cost $\mathcal C\le S_L^2/v+\sum_{\ell=0}^L\tau_\ell$.
\end{enumerate}
\end{proposition}

\begin{proof}
(i) By \cref{def:ml-coupled}, $\E[Y_\ell]=\E[P_\ell]-\E[P_{\ell-1}]$ for $\ell\ge1$; only the
two means enter. Each $\hat Y_\ell$ is unbiased for $\E[Y_\ell]$, and the sum telescopes as
in \eqref{eq:ml-telescoping}.

(ii) The $\hat Y_\ell$ are built from disjoint sets of independent samples, so they are
independent and their variances add; $\Var(\hat Y_\ell)=V_\ell/n_\ell$ by
\cref{prop:mc-unbiased}(iii). The mean squared error is the squared bias plus the variance,
as in \cref{prop:mc-mse}(i).

(iii) Apply the Cauchy--Schwarz inequality to $a_\ell=\sqrt{V_\ell/n_\ell}$ and
$b_\ell=\sqrt{n_\ell\tau_\ell}$: $(\sum a_\ell^2)(\sum b_\ell^2)\ge(\sum a_\ell b_\ell)^2=S_L^2$,
with equality if and only if $a=tb$ for some $t\ge0$ (as $b\ne0$). If every $V_\ell>0$, then
$t>0$, and $\sqrt{V_\ell/n_\ell}=t\sqrt{n_\ell\tau_\ell}$ means $n_\ell=t^{-1}\sqrt{V_\ell/\tau_\ell}$.

(iv) The lower bound is (iii) with $\Var(\hat\theta_L)\le v$. For \eqref{eq:ml-nl},
$n_\ell\ge v^{-1}S_L\sqrt{V_\ell/\tau_\ell}$ gives $V_\ell/n_\ell\le v\sqrt{V_\ell\tau_\ell}/S_L$
(trivially if $V_\ell=0$), and summing, $\Var(\hat\theta_L)\le v$. And
$n_\ell\le v^{-1}S_L\sqrt{V_\ell/\tau_\ell}+1$, so
$\sum n_\ell\tau_\ell\le v^{-1}S_L\sum\sqrt{V_\ell\tau_\ell}+\sum\tau_\ell$.
\end{proof}

Plain Monte Carlo on the finest level reaches the variance $v$ at the cost
$\Var(P_L)\,2^L/v$, one fine path of $2^L$ steps per sample. The multilevel estimator replaces the
work-normalised variance $\Var(P_L)\,2^L$ (\cref{def:mc-efficiency}) by $S_L^2$, which is much
smaller when $V_\ell\ll\Var(P_\ell)$ at the expensive levels. The optimal $n_\ell$ spends the
amount $n_\ell\tau_\ell\propto\sqrt{V_\ell\tau_\ell}$ on level $\ell$; where that amount
concentrates decides the complexity (\cref{sec:ml-complexity}). Two degenerate cases are worth
checking. If $V_\ell=0$, then $Y_\ell$ equals its mean almost surely: one sample is exact, the
continuous optimum $n_\ell=0$ is not admissible because $\E[Y_\ell]$ must still be added, and
\eqref{eq:ml-nl} takes $n_\ell=1$, which is the reason for the term $\sum\tau_\ell$ in (iv);
this happens for an equation with constant coefficients, Brownian motion with drift, for which
the Euler scheme is exact on every grid, so that the fine and coarse values of $X_T$ coincide (with
additive noise but a state-dependent drift, as for Ornstein--Uhlenbeck, it is not, and $V_\ell>0$). If
$L=0$, the estimator is plain Monte Carlo with one coarse step.

\begin{remark}[A chain of control variates]\label{rem:ml-control}
With $L=1$, $\hat\theta_1=\frac1{n_1}\sum_i\big(P^{\mathrm f}_1-P^{\mathrm c}_0\big)^{(i)}+\frac1{n_0}\sum_jP_0^{(j)}$
is the control-variate estimator \eqref{eq:cv-estimator} for $P_1$, with the control
$P^{\mathrm c}_0$, the coefficient $1$, and the unknown control mean $\E[P_0]$ replaced by an
independent estimate from $n_0$ cheap samples. Any coefficient $b$ in place of $1$, used in both
sums, keeps the estimator unbiased; its variance
$\big(\Var(P_1)-2b\,\Cov(P^{\mathrm f}_1,P^{\mathrm c}_0)+b^2\Var(P_0)\big)/n_1+b^2\Var(P_0)/n_0$ is
smallest for $b^*=\Cov(P^{\mathrm f}_1,P^{\mathrm c}_0)/\big(\Var(P_0)(1+n_1/n_0)\big)$, the
coefficient of \cref{thm:cv-optimal} shrunk to account for the noise in the estimate of the control
mean. When the levels are close, the correlation of the two members is near $1$, their variances are
near $\Var(P)$, and $n_1\ll n_0$, so $b^*\approx1$: the simple choice $b=1$ loses little. A multilevel estimator
is a chain of such controls: each level controls the one above it, and its own mean is
estimated with the help of the level below. Computing both members from one Brownian path is
common random numbers (\cref{prop:cv-crn}); with independent paths,
$V_\ell=\Var(P_\ell)+\Var(P_{\ell-1})$ does not shrink and the method loses
(\exref{ex:ml-uncoupled}).
\end{remark}

\section{Coupling fine and coarse paths}\label{sec:ml-coupling}

Let $Y^h$ denote the Euler scheme \eqref{eq:ps-euler} or the Milstein scheme
\eqref{eq:ps-milstein} with step $h$, driven by the increments of $W$ on its grid. One sample of
$Y_\ell$, $\ell\ge1$, uses one draw of the $2^\ell$ fine increments
$\Delta W_1,\dots,\Delta W_{2^\ell}$ of $W$ over steps of length $h_\ell$: the fine member is
the payoff of $Y^{h_\ell}$ driven by them, the coarse member the payoff of $Y^{h_{\ell-1}}$ driven
by the sums $\Delta W_{2k-1}+\Delta W_{2k}$, $k\le2^{\ell-1}$, which are the increments of the same
$W$ over the coarse grid (\texttt{paths.coarsen}). Both members approximate the payoff of the
same $X$. The sums are independent $\mathcal N(0,h_{\ell-1})$ variables, so the coarse member
is computed by the level-$(\ell-1)$ scheme from inputs with the right law: it has the law of
$P_{\ell-1}$, and $Y_\ell$ is a coupled correction. (Drawing the coarse increments first and
splitting each one with the Brownian bridge of \cref{prop:ps-bridge} gives the same joint law.)

\begin{proposition}[Variance of the corrections: Euler against Milstein]\label{prop:ml-variance-decay}
Let $Y_\ell$ be built as above from the Euler or the Milstein scheme, and suppose that for some
$p\ge1$ there are constants $A_p,\lambda>0$ with
\begin{equation}\label{eq:ml-strong}
  \big(\E|X_T-Y^h_T|^p\big)^{1/p}\le A_p\,h^\lambda\qquad\text{for all }h=T2^{-\ell},\ \ell\in\N .
\end{equation}
\begin{enumerate}[label=(\roman*)]
  \item (Lipschitz payoffs.) Let $P=f(X_T)$ and $P_\ell=f(Y^{h_\ell}_T)$ with
    $|f(x)-f(z)|\le\Lambda|x-z|$. If \eqref{eq:ml-strong} holds with $p=2$, then
    \[
      V_\ell\le\E[Y_\ell^2]\le\Lambda^2A_2^2\,(1+2^\lambda)^2\,h_\ell^{2\lambda}.
    \]
  \item (Digital payoffs.) Let $P=\1_{\{X_T>K\}}$ and $P_\ell=\1_{\{Y^{h_\ell}_T>K\}}$. If $X_T$
    has a density bounded by $\rho$ on $[K-\delta_0,K+\delta_0]$ for some $\delta_0>0$, then for
    every $\ell$ with $h_\ell^{\lambda p/(p+1)}\le\delta_0$,
    \[
      V_\ell\le\E[Y_\ell^2]=\Q\big(\1_{\{Y^{h_\ell}_T>K\}}\ne\1_{\{Y^{h_{\ell-1}}_T>K\}}\big)
      \le\big(2\rho+(1+2^{\lambda p})A_p^p\big)\,h_\ell^{\lambda p/(p+1)} .
    \]
  \item For the scalar equation \eqref{eq:ps-sde}, the bound \eqref{eq:ml-strong} holds with $p=2$,
    $\lambda=\frac12$ for the Euler scheme under the hypotheses (A1)--(A3) of \cref{thm:ps-euler},
    and with $p=2$, $\lambda=1$ for the Milstein scheme under the hypotheses of
    \cref{thm:ps-milstein}. Consequently, for Lipschitz payoffs $V_\ell=O(h_\ell)$ with Euler and
    $V_\ell=O(h_\ell^2)$ with Milstein. (For a system driven by several Brownian motions with
    non-commuting diffusion vector fields, the Milstein scheme without the L\'evy areas has strong
    order only $\frac12$, and then $\beta=1$ \citep{giles-szpruch}.)
\end{enumerate}
\end{proposition}

\begin{proof}
(i) $|Y_\ell|\le\Lambda|Y^{h_\ell}_T-Y^{h_{\ell-1}}_T|\le\Lambda\big(|Y^{h_\ell}_T-X_T|+|X_T-Y^{h_{\ell-1}}_T|\big)$.
Both schemes are driven by the same $W$ as $X$, so \eqref{eq:ml-strong} applies to each
term, and Minkowski's inequality gives
$\|Y_\ell\|_2\le\Lambda A_2(h_\ell^\lambda+h_{\ell-1}^\lambda)=\Lambda A_2(1+2^\lambda)h_\ell^\lambda$.
Finally $V_\ell\le\E[Y_\ell^2]$.

(ii) $Y_\ell$ takes values in $\{-1,0,1\}$, so $\E[Y^2_\ell]$ is the probability that the two
indicators differ. Let $0<\delta\le\delta_0$. If $|X_T-K|>\delta$ and both schemes are within
$\delta$ of $X_T$, then both lie strictly on the side of $K$ where $X_T$ lies, and the
indicators agree. Hence, by Markov's inequality and \eqref{eq:ml-strong},
\[
  \E[Y^2_\ell]\le\Q(|X_T-K|\le\delta)+\frac{\E|Y^{h_\ell}_T-X_T|^p+\E|Y^{h_{\ell-1}}_T-X_T|^p}{\delta^p}
  \le2\rho\delta+(1+2^{\lambda p})A_p^p\,\frac{h_\ell^{\lambda p}}{\delta^p}.
\]
The choice $\delta=h_\ell^{\lambda p/(p+1)}$ balances the two terms and gives the bound.

(iii) \Citet{kloeden-platen} prove \cref{thm:ps-euler}(i) and \cref{thm:ps-milstein}(i) by
bounding $\E|X_T-Y^h_T|^2$ (Theorems~10.2.2 and 10.3.5 there; see also the comment after
the proof of \cref{prop:ps-richardson}), which is \eqref{eq:ml-strong} with $p=2$; then apply (i).
\end{proof}

So the variance of the corrections decays at the rate $\beta=2\lambda$, twice the strong order:
$\beta=1$ for Euler, $\beta=2$ for Milstein. This is where the strong order, irrelevant for the
bias of prices (\cref{sec:ps-schemes}), pays off. The coupling is essential to the proof: the
two schemes are compared with the same $X_T$.

For digital payoffs the variance decays more slowly, because a small difference between the fine
and coarse values changes the payoff by a full unit whenever $X_T$ is near the strike. With $p=2$,
(ii) gives $\beta\ge\frac13$ for Euler and $\beta\ge\frac23$ for Milstein. For an autonomous
equation with Lipschitz coefficients, such as GBM, the Euler bound \eqref{eq:ml-strong} holds for
every $p\ge2$
(\citealp{giles-higham-mao}, equation~(1.8), referring to \citealp{kloeden-platen}), and letting
$p\to\infty$ in (ii) gives $V_\ell=O(h_\ell^{1/2-\delta})$ for every $\delta>0$; this is the
argument of \textcite{giles-higham-mao}, Section~3, and \textcite{avikainen} proves the sharp form
of the inequality behind it (Theorem~3.4 of the preprint arXiv:0712.3635). \Citet{giles-acta},
Table~5.2, lists $O(h^{1/2}\log h)$ for the Euler digital. One cannot hope for much more with Euler:
for GBM, the supremum over strikes $K\ge K_0$ of the strong error of the binary payoff is of exact
order $h^{1/2}$ (Theorem~7.2 of the preprint), whereas corrections with $\E[Y_\ell^2]\le Ch_\ell^\beta$
uniformly in $K\ge K_0$ for some $\beta>\frac12$ would give
$\E|P-P_L|=\|P-P_L\|_2^2\le\big(\sum_{\ell>L}\|Y_\ell\|_2\big)^2=O(h_L^\beta)$. For Milstein,
\textcite{avikainen} proves that the strong error of the binary payoff is $O(h^{1-\delta})$ for every
$\delta>0$ (scalar equation with Lipschitz coefficients and Lipschitz first derivatives in $x$, and
$X_T$ with a bounded density; Theorem~3.7 of the preprint). Since
$\E[Y_\ell^2]\le\E|P_\ell-P|+\E|P_{\ell-1}-P|$ for indicators, $V_\ell=O(h_\ell^{1-\delta})$;
\textcite{giles-debrabant-rossler} analyse the Milstein corrections in detail.

\begin{example}[Exact level variances for geometric Brownian motion]\label{exa:ml-gbm-exact}
Take Euler for GBM and the payoff $P=DS_T$, $D=e^{-rT}$, whose level variances can be computed
exactly. Group the fine steps in pairs. Over one coarse step the fine path is multiplied by
$f_1f_2$ with $f_i=a+\sigma\Delta W_i$, $a=1+\mu h$, $h=h_\ell$, and the coarse path by
$F=A+\sigma(\Delta W_1+\Delta W_2)$, $A=1+2\mu h$. The $n=2^{\ell-1}$ coarse steps contribute
independent factors, so with $u=\E[(f_1f_2)^2]=(a^2+\sigma^2h)^2$,
$w=\E[f_1f_2F]=a^2A+2a\sigma^2h$ and $v=\E[F^2]=A^2+2\sigma^2h$,
\[
  V_\ell=D^2S_0^2\Big[u^n-2w^n+v^n-\big(a^{2n}-A^n\big)^2\Big].
\]
Here $u-2w+v=\E[(f_1f_2-F)^2]=\sigma^4h^2+2\mu^2\sigma^2h^3+\mu^4h^4$ and
$w-v=(A\mu^2+2\mu\sigma^2)h^2$. Writing $u^n-2w^n+v^n=v^n\big[(1+x)^n-2(1+y)^n+1\big]$ with
$x=(u-v)/v$, $y=(w-v)/v$, both $O(h^2)$, and expanding the powers with $n=T/(2h)$ gives
$n(x-2y)+O(n^2h^4)=\frac T2\sigma^4h+O(h^2)$, while $(a^{2n}-A^n)^2=O(h^2)$ and
$v^n\to e^{(2\mu+\sigma^2)T}$. Hence
\[
  V_\ell=\tfrac12\,\sigma^4T\,\E[P^2]\,h_\ell+O(h_\ell^2),\qquad \E[P^2]=D^2S_0^2e^{(2\mu+\sigma^2)T}.
\]
For the test case $\E[P^2]=10^4$ exactly, so $V_\ell=8h_\ell+O(h_\ell^2)$: $V_1=3.9306$,
$V_4=0.49890$, $V_7=0.062483$. The rate $\beta=1$ of \cref{prop:ml-variance-decay} is attained,
not just a bound. For Milstein the same computation gives $V_\ell$ of exact order $h_\ell^2$,
with a constant that receives a second contribution (\exref{ex:ml-milstein-exact}).
\end{example}

\begin{figure}[ht]
  \centering
  \includegraphics[width=\textwidth]{ch12-levels}
  \caption{Level diagnostics for the test case: Euler and Milstein for GBM, the discounted
    at-the-money call and digital, levels $\ell=0,\dots,10$ (call) and $0,\dots,7$ (digital);
    $5\times10^5$ coupled paths per level for $\ell\le7$, on which both payoffs are evaluated,
    and $4\times10^4$ for $\ell=8,9,10$. Our implementation (NumPy \texttt{default\_rng}, seed
    $[2027,s,\ell]$ for scheme $s$ and level $\ell$). Left: $\log_2V_\ell$ for $\ell\ge1$, and
    $\log_2\Var(P_\ell)$ for Milstein (grey). Right: $\log_2$ of the absolute sample mean of
    $Y_\ell$, $\ell\ge1$; the Milstein digital is omitted, its means being below their
    standard errors from $\ell=2$ on.}
  \label{fig:ml-levels}
\end{figure}

\Cref{fig:ml-levels} shows these rates on the test case. For the call, $V_\ell$ halves at
each level with Euler (fitted slope $0.99$ for $\ell\ge4$; $V_{10}=4.6\times10^{-3}$) and is
divided by four with Milstein (slope $1.99$; $V_{10}=2.0\times10^{-7}$), while $\Var(P_\ell)$
stays near $158$: at level $10$ a correction has a variance $3\times10^4$ (Euler) and almost
$10^9$ (Milstein) times smaller than the payoff. The means fall like $2^{-\ell}$ for both
schemes (fitted slope $1.06$ for Milstein, $\ell\ge2$; the Euler means at $\ell\ge8$ are within
two standard errors of $0$), the weak order $\alpha=1$ of \cref{thm:ps-euler,thm:ps-milstein}. For the
digital, Euler's $V_\ell$ falls only by $\sqrt2$ per level (fitted slope $0.50$ for $\ell\ge4$), and
Milstein's by factors between $1.5$ and $2.3$ for $\ell\ge4$ (fitted slope $0.85$, resting on about
$36$ nonzero corrections out of $5\times10^5$ at $\ell=7$). A digital correction is almost always $0$, and its kurtosis
$\E[(Y_\ell-\E Y_\ell)^4]/V_\ell^2$ is large: $260$ for Euler and $1.4\times10^4$ for
Milstein at $\ell=7$, against $7.6$ for the Euler call. This matters for estimating $V_\ell$
(\cref{rem:ml-kurtosis}).

\demo{12-multilevel.ipynb}{level variances and means for the call, the digital and an Asian
option, with Euler or Milstein, for any volatility.}

\subsection*{Path-dependent payoffs}

For a payoff that depends on the path, each level needs a rule that turns a discrete path into a
payoff, and the telescoping requires only that the coarse member at level $\ell$ have the mean of
the fine member at level $\ell-1$. The simplest choice applies the same rule on every grid.

\emph{Asian options.} For $A=\frac1T\int_0^TS_t\,\dd t$, use the trapezoidal rule on the grid of
each path. The payoff is Lipschitz in the path values, but the coarse and fine averages also
differ by the difference of the two quadrature errors, whose variance is of order $h_\ell^2$.
For the discounted test-case Asian call $(A-K)^+$ (our implementation, seeds $[4,\ell]$,
$\max(4\times10^5\cdot2^{-\ell},2\times10^4)$ paths), the ratio $V_\ell/V_{\ell+1}$ stays between
$3.8$ and $4.2$ for $1\le\ell\le7$ with Milstein. With Euler the quadrature term dominates at
coarse levels (ratios $3.2$ to $3.4$ for $\ell\le3$), and the asymptotic ratio $2$ of the strong
order $\frac12$ emerges only slowly ($2.3$ and $2.2$ for $\ell=6,7$).
These are the rates $\beta=1$ and $2$ listed by \textcite{giles-acta}, Section~5.

\emph{Extremes and barriers.} A maximum, a lookback or a barrier inspected only on the grid has a
bias of order $h^{1/2}$ (\cref{prop:bar-bgk} for GBM, \cref{rem:bar-general} for Euler on general diffusions), so $\alpha=\frac12$, and a knock-out indicator is
discontinuous, so that with Euler $V_\ell=O(h_\ell^{1/2-\delta})$ for barriers and
$O(h_\ell^{1-\delta})$ for lookbacks \citep{giles-higham-mao}. The cure is to look between grid
points with the Brownian bridge: use the probability that the bridge over each step crosses the
barrier (\cref{prop:bar-bridge-crossing}), or sample the extreme of the bridge. This restores
$\alpha=1$, and with Milstein it gives $\beta\approx\frac32$ for barriers and $\beta\approx2$ for
lookbacks \citep{giles-milstein}. There, the coarse member uses the value of the fine path at the
midpoint of each coarse step, which changes its law but not its mean, so the identity
\eqref{eq:ml-telescoping} is preserved; this freedom is why \cref{def:ml-coupled} constrains only
means.

\section{Giles' complexity theorem}\label{sec:ml-complexity}

\begin{theorem}[Giles]\label{thm:ml-complexity}
Let $P\in L^2$ and let $Y_\ell$ be coupled corrections with variances $V_\ell$ and expected costs
$\tau_\ell>0$. Suppose that there are constants $\alpha,\beta,\gamma,c_1,c_2,c_3>0$ with
\[
  \alpha\ge\tfrac12\min(\beta,\gamma)
\]
such that for every $\ell\in\N$
\[
  \text{(B)}\ \ \big|\E[P_\ell]-\E[P]\big|\le c_12^{-\alpha\ell},\qquad
  \text{(V)}\ \ V_\ell\le c_22^{-\beta\ell},\qquad
  \text{(C)}\ \ \tau_\ell\le c_32^{\gamma\ell}.
\]
\begin{enumerate}[label=(\roman*)]
  \item There is a constant $c_4$, depending only on $\alpha,\beta,\gamma,c_1,c_2,c_3$, such that
    for every $\varepsilon\in(0,e^{-1})$ the finest level
    $L=\max\big(0,\big\lceil\alpha^{-1}\log_2(\sqrt2\,c_1/\varepsilon)\big\rceil\big)$ and the
    sample sizes \eqref{eq:ml-nl} with $v=\varepsilon^2/2$ give
    $\E[(\hat\theta_L-\theta)^2]\le\varepsilon^2$ at an expected cost
    \[
      \mathcal C\le\begin{cases}
        c_4\,\varepsilon^{-2}, & \beta>\gamma,\\
        c_4\,\varepsilon^{-2}(\log\varepsilon)^2, & \beta=\gamma,\\
        c_4\,\varepsilon^{-2-(\gamma-\beta)/\alpha}, & \beta<\gamma.
      \end{cases}
    \]
  \item Conversely, suppose that $|\E[P_\ell]-\E[P]|\ge c_1'2^{-\alpha\ell}$,
    $V_\ell\ge c_2'2^{-\beta\ell}$ and $\tau_\ell\ge c_3'2^{\gamma\ell}$ for every $\ell$, with
    constants $c_i'>0$. Then there are $c_4'>0$ and $\varepsilon_0>0$ such that, for
    $\varepsilon<\varepsilon_0$, every estimator \eqref{eq:ml-estimator} (any $L$, any
    $n_\ell\ge1$) with mean squared error at most $\varepsilon^2$ costs at least
    $c_4'\varepsilon^{-2}$, $c_4'\varepsilon^{-2}(\log\varepsilon)^2$ or
    $c_4'\varepsilon^{-2-(\gamma-\beta)/\alpha}$ in the three cases.
\end{enumerate}
\end{theorem}

\begin{proof}
(i) \emph{Bias.} If $L=0$ then $\sqrt2c_1\le\varepsilon$; otherwise
$L\ge\alpha^{-1}\log_2(\sqrt2c_1/\varepsilon)$. In both cases (B) gives
$|\E[P_L]-\theta|\le c_12^{-\alpha L}\le\varepsilon/\sqrt2$.

\emph{Variance and cost.} By \cref{prop:ml-telescoping}(iv) with $v=\varepsilon^2/2$, the
variance is at most $\varepsilon^2/2$, so the mean squared error is at most $\varepsilon^2$ by
\cref{prop:ml-telescoping}(ii), and
\begin{equation}\label{eq:ml-cost-bound}
  \mathcal C\le2\varepsilon^{-2}S_L^2+\sum_{\ell=0}^L\tau_\ell,\qquad
  S_L\le\sqrt{c_2c_3}\,\sum_{\ell=0}^L2^{(\gamma-\beta)\ell/2}
\end{equation}
by (V) and (C). Two consequences of the choice of $L$ and of $\varepsilon<e^{-1}<1$ will be used.
First, $L\le1+\alpha^{-1}\log_2^+(\sqrt2c_1/\varepsilon)$, so with $c_5=\max(1,\sqrt2c_1)^{1/\alpha}$,
\[
  2^L\le2c_5\,\varepsilon^{-1/\alpha}.
\]
Second, $\log(1/\varepsilon)>1$, so $L+1\le2+\alpha^{-1}\log_2c_5^\alpha+\alpha^{-1}\log_2(1/\varepsilon)\le c_6\log(1/\varepsilon)$
for a constant $c_6$. The last sum in \eqref{eq:ml-cost-bound}, which comes from the integer
ceilings, is at most $c_32^{\gamma L}/(1-2^{-\gamma})\le c_7\varepsilon^{-\gamma/\alpha}$.

\emph{Case $\beta>\gamma$.} The geometric series gives $S_L\le\sqrt{c_2c_3}/(1-2^{-(\beta-\gamma)/2})$,
so $2\varepsilon^{-2}S_L^2\le c\,\varepsilon^{-2}$. Here $\min(\beta,\gamma)=\gamma$, so
$\gamma/\alpha\le2$ and $\varepsilon^{-\gamma/\alpha}\le\varepsilon^{-2}$.

\emph{Case $\beta=\gamma$.} Every term of $S_L$ equals $\sqrt{c_2c_3}$, so
$S_L\le\sqrt{c_2c_3}\,(L+1)\le\sqrt{c_2c_3}\,c_6\log(1/\varepsilon)$ and
$2\varepsilon^{-2}S_L^2\le c\,\varepsilon^{-2}(\log\varepsilon)^2$. Again $\gamma/\alpha\le2$, and
$\varepsilon^{-2}\le\varepsilon^{-2}(\log\varepsilon)^2$.

\emph{Case $\beta<\gamma$.} Now the series is dominated by its last term:
$S_L\le\sqrt{c_2c_3}\,2^{(\gamma-\beta)L/2}/(1-2^{-(\gamma-\beta)/2})$, and with
$2^L\le2c_5\varepsilon^{-1/\alpha}$, $2\varepsilon^{-2}S_L^2\le c\,\varepsilon^{-2-(\gamma-\beta)/\alpha}$.
Here $\min(\beta,\gamma)=\beta\le2\alpha$, which is the same as
$\gamma/\alpha\le2+(\gamma-\beta)/\alpha$, so the ceiling term is of no larger order.

(ii) The mean squared error is at least the squared bias $(\E[P_L]-\theta)^2\ge c_1'^24^{-\alpha L}$,
so $2^{-\alpha L}\le\varepsilon/c_1'$, that is $L\ge\alpha^{-1}\log_2(c_1'/\varepsilon)$. It is
also at least the variance, so $\Var(\hat\theta_L)\le\varepsilon^2$, and by
\cref{prop:ml-telescoping}(iv) $\mathcal C\ge\varepsilon^{-2}S_L^2$ with
$S_L\ge\sqrt{c_2'c_3'}\sum_{\ell\le L}2^{(\gamma-\beta)\ell/2}$. If $\beta>\gamma$, keep the
term $\ell=0$: $\mathcal C\ge c_2'c_3'\varepsilon^{-2}$. If $\beta=\gamma$,
$S_L\ge\sqrt{c_2'c_3'}(L+1)\ge\sqrt{c_2'c_3'}\,\alpha^{-1}\log_2(c_1'/\varepsilon)$, which is at
least a constant times $\log(1/\varepsilon)$ for small $\varepsilon$. If $\beta<\gamma$, keep the
term $\ell=L$: $S_L^2\ge c_2'c_3'2^{(\gamma-\beta)L}\ge c_2'c_3'(c_1'/\varepsilon)^{(\gamma-\beta)/\alpha}$.
\end{proof}

This is Theorem~3.1 of \textcite{giles-mlmc}, stated there with $\gamma=1$, $\alpha\ge\frac12$ and
a general refinement factor $M$ ($h_\ell=TM^{-\ell}$, which only changes the constants); the
present form is Theorem~2.1 of \textcite{cliffe-giles-scheichl-teckentrup}, proved in their
Appendix~A, and \textcite{giles-acta}, Section~2.1 (Theorem~1 in the author's version). Part (ii) shows that the three rates are not artefacts of the
proof: for a hierarchy whose bias, variances and costs behave exactly like powers of $2$, no
choice of $L$ and $n_\ell$ does better. Where the work goes explains the three regimes. With
$n_\ell\tau_\ell\propto\sqrt{V_\ell\tau_\ell}\propto2^{(\gamma-\beta)\ell/2}$, the work
concentrates on the coarsest levels if $\beta>\gamma$, so the cost is that of an unbiased
estimator with bounded cost per sample; it is spread evenly over the $L+1\approx\alpha^{-1}\log_2(1/\varepsilon)$
levels if $\beta=\gamma$, hence the squared logarithm; and it concentrates on the finest levels if
$\beta<\gamma$.

\begin{remark}[The condition $\alpha\ge\frac12\min(\beta,\gamma)$]\label{rem:ml-alpha}
It is used only to bound the term $\sum_\ell\tau_\ell$ in \eqref{eq:ml-cost-bound}: every level up
to $L$ must be sampled at least once, and the bias fixes how many levels there are. It cannot be
dropped. Under the lower bounds of (ii), any estimator with mean squared error $\varepsilon^2$
costs at least $\sum_{\ell\le L}\tau_\ell\ge c_3'2^{\gamma L}\ge c_3'(c_1'/\varepsilon)^{\gamma/\alpha}$,
and if $\alpha<\frac12\min(\beta,\gamma)$ this exceeds the bound in (i) by a power of
$1/\varepsilon$. For the Euler and Milstein couplings with smooth or Lipschitz payoffs,
$\alpha=1$ and $\gamma=1$, and the condition holds. For a lookback monitored on the grid with
Euler, $\alpha=\frac12$ and $\beta=1-\delta$ (\cref{sec:ml-coupling}), and it holds only just.
\end{remark}

For plain Monte Carlo with a bias $c\,h^\alpha$ and a cost $h^{-\gamma}$ per path,
\cref{prop:mc-mse}(iii) gives a cost of order $\varepsilon^{-2-\gamma/\alpha}$. The table uses the
rates $\beta$ of \cref{sec:ml-coupling}, $\gamma=1$, and for $\alpha$ the weak orders of
\cref{thm:ps-euler,thm:ps-milstein} (for the digital, those of the expansion of \citealp{bally-talay}
recalled in \cref{rem:ps-failures}, also observed in \cref{fig:ml-levels}); the rate for the
Milstein digital is Avikainen's, recalled after \cref{prop:ml-variance-decay} (a sample of $Y_\ell$ costs $2^\ell+2^{\ell-1}=3\cdot2^{\ell-1}$
steps):
\begin{center}
\begin{tabular}{lcccc}
\toprule
payoff and scheme & $\alpha$ & $\beta$ & multilevel & plain Monte Carlo\\
\midrule
Lipschitz, Euler & $1$ & $1$ & $\varepsilon^{-2}(\log\varepsilon)^2$ & $\varepsilon^{-3}$\\
Lipschitz, Milstein & $1$ & $2$ & $\varepsilon^{-2}$ & $\varepsilon^{-3}$\\
digital, Euler & $1$ & $\frac12$ & $\varepsilon^{-5/2}$ & $\varepsilon^{-3}$\\
digital, Milstein & $1$ & $1-\delta$ & $\varepsilon^{-2-\delta}$ & $\varepsilon^{-3}$\\
barrier on the grid, Euler & $\frac12$ & $\frac12-\delta$ & $\varepsilon^{-3-2\delta}$ & $\varepsilon^{-4}$\\
\bottomrule
\end{tabular}
\end{center}
For the digital the slower decay of $V_\ell$ moves Euler from the borderline regime $\beta=\gamma$
into the regime $\beta<\gamma$: the cost of the finest levels dominates, and halving $\varepsilon$
multiplies the cost by $2^{5/2}\approx5.7$ instead of about $4$. Milstein restores $\beta\approx\gamma$,
and smoothing the payoff by a conditional expectation over the last step pushes $\beta$ above
$\gamma$ (\exref{ex:ml-smoothing}; \citealp{giles-acta}, Section~5).

\subsection*{Against the optimal allocation of Chapter 7}

\Cref{prop:ps-allocation} chose the best single step for a given budget. \Cref{prop:ml-telescoping}
chooses the best spread of the same budget over a hierarchy of steps. On the laboratory of
\cref{fig:ps-allocation} (GBM with $\mu=-1\%$, $\sigma=20\%$, the smooth payoff
$f(x)=5\log(1+e^{(x-100)/5})$, cost in steps), the multilevel root-mean-square error for a budget
$\mathcal C$ is $\min_L\big(b_L^2+S_L^2/\mathcal C\big)^{1/2}$ (continuous sample sizes, as in
\cref{prop:ml-telescoping}(iii)), with the exact biases $b_L$
(\texttt{paths.gbm\_scheme\_weak\_error}) and the level variances estimated with
$\max(4\times10^5\cdot2^{-\ell},2\times10^4)$ paths (our implementation, seeds $[3,\ell]$ for
Euler and $[4,\ell]$ for Milstein; finest levels $L=4,8$ for Euler and $6,10$ for Milstein). For a
fair comparison, the single-level columns use the same exact biases: the error is
$\min_m\big(b_m^2+v\,m/\mathcal C\big)^{1/2}$ over integer numbers of steps $m$, with $v=154$ (best
$m=7,34$ for Euler and $5,19$ for Milstein; the model of \cref{prop:ps-allocation} with $c=0.17$
gives $0.041$ and $0.0088$ for Euler):
\begin{center}
\begin{tabular}{lcccccc}
\toprule
budget (steps) & exact & Euler & Milstein & MLMC, Euler & MLMC, Milstein\\
\midrule
$10^6$ & $0.0124$ & $0.0400$ & $0.0318$ & $0.0227$ & $0.0136$\\
$10^8$ & $0.00124$ & $0.00885$ & $0.00659$ & $0.00307$ & $0.00137$\\
\bottomrule
\end{tabular}
\end{center}
Forty independent replications of each multilevel estimator at $10^6$ steps have root-mean-square
errors $0.026$ (Euler) and $0.012$ (Milstein) (\cref{fig:ml-complexity}, right). At equal budget,
the multilevel estimators cut the error of the best single-level estimator with the same scheme by
factors of $1.8$ and $2.9$ (Euler) and $2.3$ and $4.8$ (Milstein) at $10^6$ and $10^8$ steps, and
the error of multilevel Milstein is only about $10\%$ above that of exact simulation: in terms of
cost, it pays a factor of about $1.2$ for not knowing the law of $S_T$.

\section{The adaptive algorithm}\label{sec:ml-algorithm}

In practice $V_\ell$ and the bias are unknown, and the number of levels is decided on the fly.
The bias test rests on the following observation.

\begin{lemma}[Estimating the remaining bias]\label{lem:ml-bias}
If $\E[P_\ell]-\E[P]=c\,2^{-\alpha\ell}(1+\eta_\ell)$ with $c\ne0$, $\alpha>0$ and $\eta_\ell\to0$,
then
\[
  \E[P_L]-\E[P]=-\frac{\E[Y_L]}{2^\alpha-1}\,\big(1+o(1)\big)\qquad(L\to\infty).
\]
\end{lemma}

\begin{proof}
$\E[Y_L]=c2^{-\alpha L}\big(1+\eta_L-2^\alpha(1+\eta_{L-1})\big)=-c2^{-\alpha L}(2^\alpha-1)(1+o(1))$,
and $\E[P_L]-\E[P]=c2^{-\alpha L}(1+o(1))$; divide.
\end{proof}

\medskip\noindent\textbf{Algorithm} (\citealp{giles-mlmc}, Section~5; \citealp{giles-acta}, Section~3.1).
Given $\varepsilon$, rates $\alpha$ and $\beta$ (known from the theory, or fitted to the level
means and variances), and an initial sample size $n_{\mathrm{init}}$:
\begin{enumerate}[label=\arabic*.]
  \item Set $L=2$ and draw $n_{\mathrm{init}}$ samples on levels $0,1,2$.
  \item Estimate $V_\ell$ by the sample variance of the samples drawn so far on each level, and set
    $n_\ell$ by \eqref{eq:ml-nl} with $v=\varepsilon^2/2$; draw the missing samples. Repeat until no
    level needs more than $1\%$ additional samples.
  \item Estimate the remaining bias by $\max_{j=0,1,2}\,2^{-j\alpha}|\hat Y_{L-j}|/(2^\alpha-1)$
    (levels $L-j\ge1$). If it is at most $\varepsilon/\sqrt2$, stop and return
    $\hat\theta_L$ with the standard error $(\sum_\ell\hat V_\ell/n_\ell)^{1/2}$. Otherwise set
    $L\leftarrow L+1$, initialise $\hat V_{L}=\hat V_{L-1}2^{-\beta}$, and go to step~2.
\end{enumerate}
Step~3 applies \cref{lem:ml-bias} to the sample means, extrapolating from the two levels below
$L$ as well, so that a correction whose mean happens to be close to zero does not stop the
algorithm too early. The algorithm is a heuristic, and \textcite{giles-acta} says so: the lemma is
asymptotic, the regime it assumes may not have been reached, and the sample means are noisy, so
nothing guarantees a mean squared error below $\varepsilon^2$. Its standard error ignores the
bias, which the stopping rule only estimates. Our implementation (\texttt{multilevel.mlmc}) also
prevents $\hat V_\ell$ from falling below half of $\hat V_{\ell-1}2^{-\beta}$ for $\ell\ge2$.

\begin{remark}[Kurtosis and the variance estimates]\label{rem:ml-kurtosis}
The sample variance of $n$ \iid\ copies of $Y$ has a relative standard deviation of about
$\sqrt{(\kappa-1)/n}$, where $\kappa$ is the kurtosis of $Y$: twice the relative error of the sample
standard deviation computed in \cref{exa:mc-se-accuracy}. For a
digital correction, $Y_\ell$ is nonzero with a small probability $q_\ell$ and $\kappa\approx1/q_\ell$,
so the variance estimate is reliable only when $n_\ell q_\ell$, the number of nonzero corrections
seen, is large. With $10^3$ initial samples and $\kappa=1.4\times10^4$ (\cref{fig:ml-levels}),
the estimate of $V_\ell$ is often exactly $0$: the floor in the implementation, or an initial sample
size of order $\kappa$, is needed. A smoothed payoff has a much smaller kurtosis.
\end{remark}

\begin{figure}[ht]
  \centering
  \includegraphics[width=\textwidth]{ch12-complexity}
  \caption{Left: the adaptive algorithm on the test-case call ($\alpha=1$, $\beta=1$ for Euler and
    $2$ for Milstein, $n_{\mathrm{init}}=1000$, seed $11$; one run per accuracy): $\varepsilon^2$
    times the cost in steps, for $\varepsilon=0.05$, $0.02$, $0.01$, $0.005$. Triangles: plain Monte
    Carlo on the finest level $L$ chosen by the same run, with $\lceil2\Var(P_L)/\varepsilon^2\rceil$
    paths of $2^L$ steps ($\Var(P_L)$ from \cref{fig:ml-levels}). Right: Chapter~7's laboratory
    (\cref{fig:ps-allocation}): root-mean-square error against the budget for exact simulation, Euler
    with the optimal allocation of \cref{prop:ps-allocation} ($c=0.17$, $v=154$), and the best
    multilevel estimators (lines: model of \cref{sec:ml-complexity}; dots: $40$ replications at
    $10^4$, $10^5$ and $10^6$ steps).}
  \label{fig:ml-complexity}
\end{figure}

On the test-case call (\cref{fig:ml-complexity}, left; our implementation, seed $11$), the
algorithm stops at $L=3,4,5,6$ with Euler and $L=3,4,5,5$ with Milstein for
$\varepsilon=0.05,0.02,0.01,0.005$. With Milstein, $\varepsilon^2$ times the cost stays at
$371$--$375$ steps: the regime $\beta>\gamma$, where the cost is that of an unbiased estimator. With
Euler it grows from $675$ to $1370$, the slow logarithmic growth of the regime $\beta=\gamma$. At
$\varepsilon=0.005$ the Euler run uses $(n_0,\dots,n_6)=(23497780, 1547621, 827227, 434917, 226194,
114576, 57780)$ samples, returns $7.3703$ with standard error $0.0035$, and costs
$5.5\times10^7$ steps; plain Monte Carlo on level $6$ would cost $15$ times more, and the gap
grows like $1/(\varepsilon(\log\varepsilon)^2)$. For Milstein the gap is $27$. For the digital with
Euler ($\beta=\frac12$, $n_{\mathrm{init}}=2000$), averaged over eight runs (seeds $11$ to $18$),
$\varepsilon^2\times$cost is $10.8$, $17.4$ and $32.5$ at $\varepsilon=0.004$, $0.002$, $0.001$,
with root-mean-square errors $0.0033$, $0.0022$, $0.0010$. Unlike for the call it grows, as in the
regime $\beta<\gamma$ of \cref{thm:ml-complexity}, where it grows like $\varepsilon^{-1/2}$; at these
accuracies the growth is irregular (single runs vary by a factor of $2$), because the finest level
chosen by the algorithm jumps from one run to another.

\demo{12-multilevel.ipynb}{run the adaptive algorithm on the call or an Asian option for any accuracy and scheme, see how
it spreads the samples over the levels, and compare its cost with plain Monte Carlo.}

\begin{incode}
\coderef{code/src/mfe230d/multilevel.py}: \texttt{GBMLevels(payoff, scheme)} returns, for a level
$\ell$, coupled fine and coarse payoffs computed on shared Brownian increments; \texttt{SDELevels}
does the same for any scalar equation through \texttt{paths.euler} and \texttt{paths.milstein}.
\texttt{level\_statistics} produces the diagnostics of \cref{fig:ml-levels}, \texttt{mlmc} is the
adaptive algorithm, \texttt{optimal\_samples} implements \eqref{eq:ml-nl}, \texttt{gbm\_level\_exact}
the exact moments of \cref{exa:ml-gbm-exact}, and \texttt{single\_term} and \texttt{coupled\_sum} the
estimators of \cref{sec:ml-unbiased}.
\end{incode}

\section{Unbiased randomised estimators}\label{sec:ml-unbiased}

When $\beta>\gamma$, the multilevel cost $\varepsilon^{-2}$ is that of an unbiased estimator with
bounded cost per sample. \Citet{mcleish-debiasing} and \textcite{rhee-glynn} showed that such an
estimator exists when, in addition, the bias vanishes fast enough ($2\alpha>\gamma$): draw the level
at random, so that every level, however fine, has a chance to be sampled, and reweight. There is
then no finest level and no bias.

\begin{proposition}[Single-term randomised estimator]\label{prop:ml-unbiased}
Let $Y_0,Y_1,\dots$ be coupled corrections with $\E[P_\ell]\to\E[P]$, let $m_j=\E[Y_j^2]$, and
let $J$ be an $\N$-valued random variable, independent of the corrections, with
$p_j=\Q(J=j)>0$ for every $j$. Draw $J$, then one sample of $Y_J$, and set $Z=Y_J/p_J$.
\begin{enumerate}[label=(\roman*)]
  \item $\E[Z^2]=\sum_{j\ge0}m_j/p_j$ (finite or not), and the expected cost of one sample is
    $\E[\tau_J]=\sum_jp_j\tau_j$.
  \item If $\sum_jm_j/p_j<\infty$, then $Z$ is integrable, $\E[Z]=\E[P]$, and
    $\Var(Z)=\sum_jm_j/p_j-\E[P]^2$.
  \item (Geometric randomisation.) Let $p_j=(1-2^{-r})2^{-rj}$ with $r>0$, and suppose that
    $m_j\le c\,2^{-\beta'j}$ and $\tau_j\le c\,2^{\gamma j}$. Then $\Var(Z)<\infty$ if $r<\beta'$, and
    $\E[\tau_J]<\infty$ if $r>\gamma$. If instead $m_j\ge c'2^{-\beta'j}$ and $\tau_j\ge c'2^{\gamma j}$
    with $c'>0$, then $\Var(Z)=\infty$ for $r\ge\beta'$ and $\E[\tau_J]=\infty$ for $r\le\gamma$.
    Under (B) and (V) of \cref{thm:ml-complexity}, one can take $\beta'=\min(\beta,2\alpha)$.
  \item (No free lunch.) For every distribution $(p_j)$ with all $p_j>0$,
    \[
      \E[Z^2]\,\E[\tau_J]\ \ge\ \Big(\sum_{j\ge0}\sqrt{m_j\tau_j}\Big)^2 .
    \]
    In particular, if $m_j\ge c'2^{-\beta'j}$ and $\tau_j\ge c'2^{\gamma j}$ with $\beta'\le\gamma$, then
    no randomisation gives both a finite variance and a finite expected cost.
\end{enumerate}
\end{proposition}

\begin{proof}
(i) Conditioning on $J$ and using its independence of the corrections, by Tonelli's theorem
$\E[Z^2]=\sum_jp_j\E[Y_j^2]/p_j^2=\sum_jm_j/p_j$; the cost is $\tau_j$ on $\{J=j\}$.

(ii) By the Cauchy--Schwarz inequality,
\[
  \E|Z|=\sum_j\E|Y_j|\le\sum_j\sqrt{m_j/p_j}\,\sqrt{p_j}\le\Big(\sum_j\frac{m_j}{p_j}\Big)^{1/2}\Big(\sum_jp_j\Big)^{1/2}<\infty .
\]
So $Z$ is integrable and, by Fubini's theorem and \cref{prop:ml-telescoping}(i),
$\E[Z]=\sum_j\E[Y_j]=\lim_L\sum_{j\le L}\E[Y_j]=\lim_L\E[P_L]=\E[P]$. The variance follows from (i).

(iii) $m_j/p_j\le c\,2^{(r-\beta')j}/(1-2^{-r})$ and $p_j\tau_j\le c\,2^{(\gamma-r)j}$ give
convergent geometric series when $r<\beta'$, respectively $r>\gamma$. With the lower bounds, the
terms are bounded below by $c'2^{(r-\beta')j}/(1-2^{-r})$ and $c'(1-2^{-r})2^{(\gamma-r)j}$, which
do not tend to $0$ when $r\ge\beta'$, respectively $r\le\gamma$; then $\E[Z^2]=\infty$, and
$\Var(Z)=\infty$ because $\E[P]$ is finite. Finally,
$m_j=V_j+\E[Y_j]^2$ and $|\E[Y_j]|\le|\E[P_j]-\E[P]|+|\E[P_{j-1}]-\E[P]|\le c_1(1+2^\alpha)2^{-\alpha j}$,
so $m_j\le c_22^{-\beta j}+c_1^2(1+2^\alpha)^24^{-\alpha j}\le c\,2^{-\min(\beta,2\alpha)j}$.

(iv) The Cauchy--Schwarz inequality with $a_j=\sqrt{m_j/p_j}$ and $b_j=\sqrt{p_j\tau_j}$, as in
\cref{prop:ml-telescoping}(iii). If $\beta'\le\gamma$, then $\sqrt{m_j\tau_j}\ge c'2^{(\gamma-\beta')j/2}\ge c'$,
the series diverges, and one of the two factors on the left is infinite.
\end{proof}

So a finite variance and a finite expected cost are compatible exactly when the second moments
$m_j=V_j+\E[Y_j]^2$ shrink faster than the costs grow: if $m_j\asymp2^{-\beta'j}$ and
$\tau_j\asymp2^{\gamma j}$, by (iii) and (iv) some randomisation works if and only if $\beta'>\gamma$,
and then every geometric one with $r\in(\gamma,\beta')$ does. Under (B) and (V),
$\beta'=\min(\beta,2\alpha)$, and the condition $\min(\beta,2\alpha)>\gamma$ is stronger than the
$\beta>\gamma$ of \cref{thm:ml-complexity}: with a bias $c\,2^{-\gamma\ell/2}$, variances
$2^{-\beta\ell}$, $\beta>\gamma$, and costs $2^{\gamma\ell}$, multilevel Monte Carlo costs
$\varepsilon^{-2}$, but $m_j\ge\E[Y_j]^2\asymp2^{-\gamma j}$ and (iv) rules out every randomisation.
When the condition holds, the estimator averages \iid\ samples with random costs,
and by the central limit theorem with random costs (\cref{rem:mc-budget}) its error for a budget
$\mathcal C$ is of order $(\Var(Z)\E[\tau_J]/\mathcal C)^{1/2}$: the square-root rate of exact
simulation, with no bias to estimate and a valid confidence interval
(\cref{thm:mc-clt}). Milstein with a Lipschitz payoff qualifies ($\beta=2$, $\alpha=\gamma=1$);
Euler, with $\beta=\gamma=1$, does not, and \textcite{rhee-glynn}, Section~4, analyse what can still be done
then. Equality in (iv) holds for $p_j\propto\sqrt{m_j/\tau_j}$ (\exref{ex:ml-randomisation} treats
geometric laws); \textcite{rhee-glynn}, Theorems~3 and~4, find the randomisations that minimise the
work--variance product $\Var(Z)\E[\tau_J]$.

\begin{remark}[The coupled-sum estimator]\label{rem:ml-coupled-sum}
Rhee and Glynn's main estimator uses one path for all the levels up to the random one:
$\bar Z=\sum_{j=0}^JY_j/\Q(J\ge j)$, where $Y_j=P_j-P_{j-1}$ are all computed from one Brownian
path with $2^J$ steps, coarsened level by level, and $J$ is independent of the path. If $P_j\to P$ in
$L^2$ and
$\sum_{j\ge1}\E[(P_{j-1}-P)^2]/\Q(J\ge j)<\infty$, then $\bar Z$ is square integrable and unbiased
for $\E[P]$, and its second moment is given explicitly \citep[Theorem~1]{rhee-glynn}. Level $0$
is always included, since $\Q(J\ge0)=1$, so a cheap and noisy level~$0$ is not inflated by a
factor $1/p_0$ as in the single-term estimator. The expected cost is $\sum_j2^j\Q(J\ge j)$ steps.
\end{remark}

\begin{figure}[ht]
  \centering
  \includegraphics[width=0.68\textwidth]{ch12-unbiased}
  \caption{Randomised estimators for the test-case call with Milstein, $\Q(J=j)\propto2^{-rj}$:
    variance times mean cost per sample, in steps. Line: single-term model of
    \cref{prop:ml-unbiased}(i), with $m_j$ estimated from \cref{fig:ml-levels} for $j\le10$ and
    continued geometrically with $\beta=2$, $\gamma=1$. Dots: simulations with $4\times10^5$
    samples per value of $r$ (seeds $[5,100r]$ and $[6,100r]$). Shaded: the window
    $\gamma<r<\beta$ of \cref{prop:ml-unbiased}(iii). Horizontal lines: $\varepsilon^2$ times the cost of
    the adaptive multilevel runs of \cref{fig:ml-complexity} (mean over $\varepsilon$), and the
    variance $158.0$ of exact simulation at one step per sample.}
  \label{fig:ml-unbiased}
\end{figure}

\Cref{fig:ml-unbiased} shows the window on the test case. At $r=1.5$ (our implementation, seed
$[5,150]$, $4\times10^5$ samples), the single-term estimator returns $7.4075$ with standard error
$0.0267$, at an average cost of $3.01$ steps per sample (the expected cost is $2.987$); the
coupled-sum estimator (seed $[6,150]$) returns $7.3238$ with standard error $0.0199$ at an average
cost of $3.71$ steps (expected cost $1/(1-2^{-1/2})=3.414$; see the first caveat below). The
fourteen estimates in the figure lie between $-2.4$ and $+1.7$
standard errors of $7.36429$. As $r$ decreases to $\gamma=1$, the cost of the rare fine levels
dominates and the product blows up; at the other end the variance diverges only at $r=\beta=2$, and
so slowly that the curve barely turns up, because $m_j$ falls by four per level from
$m_1\approx0.047$ while $m_0/p_0$, about $300$, dominates. Near $r=1.75$--$1.9$ the coupled sum
costs $340$--$390$ steps per unit of variance, against $373$ for $\varepsilon^2$ times the cost of
adaptive multilevel Monte Carlo, which spends half of $\varepsilon^2$ on the bias. Two caveats: for
$r\le2$ the cost per sample has infinite variance ($\E[\tau_J^2]=\sum_jp_j\tau_j^2$ diverges), so
the realised cost of a run fluctuates a lot, although the central limit theorem with random costs
needs only $\E[\tau_J]<\infty$; and the sample variance of $Z$ converges slowly when $\E[Z^4]=\infty$,
which for the single-term estimator happens when $r\ge\frac43$, since
$\E[Z^4]=\sum_j\E[Y_j^4]/p_j^3$ and $\E[Y_j^4]$ is here of the order of $m_j^2\approx4^{-2j}$.

\demo{12-multilevel.ipynb}{randomised estimators for the call with Euler or Milstein: running means,
confidence intervals and running costs for any rate $r$.}

\furtherreading
\textcite{giles-acta} is the standard survey: the complexity theorem and its variants
(Sections~2.1--2.4, including randomised and multi-index methods), the algorithm (Section~3.1), and
the analysis of path-dependent payoffs with Euler and Milstein (Section~5), with references.
\textcite{giles-mlmc} is the original paper, \textcite{giles-milstein} the Milstein version with the
Brownian-bridge treatment of lookbacks, barriers and digitals, and \textcite{giles-higham-mao} and
\textcite{avikainen} the analysis of nonsmooth payoffs with Euler. \textcite{kebaier-romberg} studied
the two-level estimator and its central limit theorem (\exref{ex:ml-kebaier}), and
\textcite{ben-alaya-kebaier} proved a central limit theorem for the multilevel Euler estimator. In
several dimensions, where Milstein needs Lévy areas, the antithetic coupling of
\textcite{giles-szpruch} recovers $\beta=2$ without them for smooth payoffs (about $\frac32$ for a
call). \textcite{rhee-glynn} and
\textcite{mcleish-debiasing} are the sources for \cref{sec:ml-unbiased}. \textcite{glasserman-mc},
which predates multilevel methods, covers the ingredients: discretisation schemes and their orders
(Sections~6.1--6.2) and the allocation of effort (Section~1.1.3).
